\documentclass[handout,t,fleqn,leqno,aspectratio=1610,9pt]{beamer}
% 11pt default
% draft, handout

\usepackage[utf8]{inputenc}
\usepackage[T1]{fontenc}
\usepackage[ngerman]{babel}
\usepackage{lmodern}
\usepackage{exscale}
\usepackage{amsmath}
\usepackage{amssymb}
\usepackage{accents}
\usepackage{textcomp}
\usepackage{gensymb}
\usepackage{csquotes}
\usepackage{dsfont}
\usepackage{microtype}
\usepackage{xspace}
\usepackage{siunitx}
\usepackage{subfigure}
\usepackage{tikz}
\usepackage{pifont}
\usepackage{listings}
\usepackage{tikzsymbols}
\usepackage[natbib=true,backend=biber,bibencoding=latin1,backref=true,style=numeric,maxnames=5,
giveninits=true,sorting=nyt]{biblatex}





\pdfsuppresswarningpagegroup=1



\graphicspath{%
{./figures/}%
}


\tikzstyle{every picture}=[remember picture, overlay]
\tikzstyle{every node}=[anchor=south west, inner sep=0pt, outer sep=0pt]



% color

% \usecolortheme{beaver}
% \usecolortheme{seahorse}
% \usecolortheme{rose}


\definecolor{checkred}   {RGB} {200,0,0}
\definecolor{checkgreen} {RGB} {123,137,84}
\definecolor{checkgold}  {RGB} {200,200,0}
\definecolor{uzkblau}    {RGB} {0,81,118}
\definecolor{uzkturkis}  {RGB} {0,157,204}
\definecolor{uzkkorall}  {RGB} {234,86,79}
\newcommand{\textc}[1]{\textcolor{uzkblau}{#1}}

% \definecolor{myblue}{RGB}{53, 126, 199}
% \definecolor{myblue}{RGB}{43, 116, 189}
% \definecolor{myblue}{RGB}{33, 106, 179}
\definecolor{myblue}{RGB}{23, 96, 169}
% \definecolor{myblue}{RGB}{3, 76, 149}
\setbeamercolor*{structure}{fg=uzkblau}

\setbeamercolor{palette primary}{use=structure,fg=black,bg=structure.fg!30!white}
\setbeamercolor{palette secondary}{use=structure,fg=black,bg=structure.fg!35!white}
\setbeamercolor{palette tertiary}{use=structure,fg=black,bg=structure.fg!40!white}
\setbeamercolor{palette quaternary}{use=structure,fg=black,bg=structure.fg!45!white}
\setbeamercolor*{titlelike}{parent=palette primary}

% \setbeamercolor{sidebar}{use=structure,bg=structure.fg!20!white}
% \setbeamercolor{palette sidebar primary}{use=normal text,fg=normal text.fg}
% \setbeamercolor{palette sidebar secondary}{use=structure,fg=structure.fg}
% \setbeamercolor{palette sidebar tertiary}{use=normal text,fg=normal text.fg}
% \setbeamercolor{palette sidebar quaternary}{use=structure,fg=structure.fg}
% \setbeamercolor*{separation line}{}
% \setbeamercolor*{fine separation line}{}



% font

\usefonttheme[onlymath]{serif}

% \setbeamerfont{framesubtitle}{size=\large}



%  bullet points

\useinnertheme{circles}

\settowidth{\leftmargini}{\usebeamertemplate{itemize item}}
\addtolength{\leftmargini}{\labelsep}





% no navigation
\beamertemplatenavigationsymbolsempty

\defbeamertemplate{footline}{toc and frame number}
{%
  \usebeamercolor[fg]{page number in head/foot}%
  \usebeamerfont{page number in head/foot}%
  \hspace{.2cm}%
  \insertauthor,
  \insertinstitute
  \hfill%
  \setbeamertemplate{page number in head/foot}[framenumber]%
  \usebeamertemplate*{page number in head/foot}\kern1em\vskip2pt%
}

\setbeamertemplate{footline}[toc and frame number]


\newenvironment{wideitemize}{\itemize\addtolength{\itemsep}{.5\baselineskip}}{\enditemize}
\newenvironment{widedescription}{\description\addtolength{\itemsep}{.5\baselineskip}}{\enddescription}


% defines

\newcommand{\citepaper}[1]{{\footnotesize \color{myblue} [#1]}}
\newcommand{\follows}{$\longrightarrow$\@\xspace}
\newcommand{\eg}{e.\,g.\@\xspace}


\newcommand{\good}{\textcolor{checkgreen}{\Large\ding{51}}}%
\newcommand{\new}{\textcolor{checkgold}{\Large\ding{72}}}%
\newcommand{\nogood}{\textcolor{checkred}{\Large\ding{56}}}%
\newcommand{\todo}{{\Large\ding{229}}}%

%\newcommand{\good}{{\color{checkgreen}\CheckmarkBold}}
%\newcommand{\nogood}{{\color{checkred}\XSolidBold}}


\lstset{
%   backgroundcolor=\color{white},   % choose the background color; you must add \usepackage{color} or \usepackage{xcolor}; should come as last argument
  basicstyle=\small,        % the size of the fonts that are used for the code
%   breakatwhitespace=false,         % sets if automatic breaks should only happen at whitespace
%   breaklines=true,                 % sets automatic line breaking
%   captionpos=b,                    % sets the caption-position to bottom
%   commentstyle=\color{mygreen},    % comment style
%   deletekeywords={...},            % if you want to delete keywords from the given language
%   escapeinside={\%*}{*)},          % if you want to add LaTeX within your code
%   extendedchars=true,              % lets you use non-ASCII characters; for 8-bits encodings only, does not work with UTF-8
%   firstnumber=1000,                % start line enumeration with line 1000
%   frame=single,	                   % adds a frame around the code
%   keepspaces=true,                 % keeps spaces in text, useful for keeping indentation of code (possibly needs columns=flexible)
  keywordstyle=\color{blue},       % keyword style
  language=C++,                 % the language of the code
  morekeywords={function, end, @kernel, @index},            % if you want to add more keywords to the set
%   numbers=left,                    % where to put the line-numbers; possible values are (none, left, right)
%   numbersep=5pt,                   % how far the line-numbers are from the code
%   numberstyle=\tiny\color{mygray}, % the style that is used for the line-numbers
%   rulecolor=\color{black},         % if not set, the frame-color may be changed on line-breaks within not-black text (e.g. comments (green here))
%   showspaces=false,                % show spaces everywhere adding particular underscores; it overrides 'showstringspaces'
%   showstringspaces=false,          % underline spaces within strings only
%   showtabs=false,                  % show tabs within strings adding particular underscores
%   stepnumber=2,                    % the step between two line-numbers. If it's 1, each line will be numbered
%   stringstyle=\color{mymauve},     % string literal style
   tabsize=2,	                   % sets default tabsize to 2 spaces
%   title=\lstname                   % show the filename of files included with \lstinputlisting; also try caption instead of title
}



% temporary grid
% \setbeamertemplate{background}[grid][step=.5cm]

% \includeonlyframes{current}



\title                              {TrixiAtmo.jl}
\author    [Atiano el al.]          {Marco Artiano, Philipp Baasch, Benedict Geihe, Tristan Montoya, Andrés Rueda-Ramírez}
\date      [11.02.26]               {TRUDI\\[6pt]March, 16\textsuperscript{th} 2026}



\begin {document}

%%%%%%%%%%%%%%%%%%%%%%%%%%%%%%%%%%%%%%%%%%%%%%%%%%%%%%%%%%%%%%%%%%%%%%%%%%%%%%%%%%%%%%%%%%
\begin{frame}
\titlepage
\end{frame}
%%%%%%%%%%%%%%%%%%%%%%%%%%%%%%%%%%%%%%%%%%%%%%%%%%%%%%%%%%%%%%%%%%%%%%%%%%%%%%%%%%%%%%%%%%



%%%%%%%%%%%%%%%%%%%%%%%%%%%%%%%%%%%%%%%%%%%%%%%%%%%%%%%%%%%%%%%%%%%%%%%%%%%%%%%%%%%%%%%%%%
\begin{frame}{Fact sheet}
% [Benedict]
\begin{wideitemize}
 \item Initial motivation: very specific requirements in projects ADAPTEX and ICON-DG \\[6pt]

 \begin{wideitemize}
    \item ``Primite equations'' (i.e. at least moisture and precipitation, \texttt{CompressibleEulerMulticomponentEquations3D} does not work)
    \item Covariant equations, PDEs on manifolds
 \end{wideitemize}

 \item Hardest part: figuring out the name \texttt{TrixiAtmo\textcolor{red}{s}.jl}

 \item Initial commit: June, 6\textsuperscript{th} 2024

 \item Official release: January, 16\textsuperscript{th} 2026

 \item 12 \new

 \item Contributors: Marco, Arpit, Daniel, Gregor, Benedict, Fabian, Oswald, Joshua, Tristan, Hendrik, Andrés, Michael, Andrew
\end{wideitemize}

\end{frame}
%%%%%%%%%%%%%%%%%%%%%%%%%%%%%%%%%%%%%%%%%%%%%%%%%%%%%%%%%%%%%%%%%%%%%%%%%%%%%%%%%%%%%%%%%%


%%%%%%%%%%%%%%%%%%%%%%%%%%%%%%%%%%%%%%%%%%%%%%%%%%%%%%%%%%%%%%%%%%%%%%%%%%%%%%%%%%%%%%%%%%
\begin{frame}{Physical background}
% [Marco]

Atmospheric flows

What makes it challenging / exciting / unique

low Mach

gravity -> stratification

rotation, Coriolis

curvature -> curvilinear, covariant

vertical direction, vastly different spatial and time scales

\end{frame}
%%%%%%%%%%%%%%%%%%%%%%%%%%%%%%%%%%%%%%%%%%%%%%%%%%%%%%%%%%%%%%%%%%%%%%%%%%%%%%%%%%%%%%%%%%


%%%%%%%%%%%%%%%%%%%%%%%%%%%%%%%%%%%%%%%%%%%%%%%%%%%%%%%%%%%%%%%%%%%%%%%%%%%%%%%%%%%%%%%%%%
\begin{frame}{Structure}
% [?]

sketch

long list of equations....

\end{frame}
%%%%%%%%%%%%%%%%%%%%%%%%%%%%%%%%%%%%%%%%%%%%%%%%%%%%%%%%%%%%%%%%%%%%%%%%%%%%%%%%%%%%%%%%%%



%%%%%%%%%%%%%%%%%%%%%%%%%%%%%%%%%%%%%%%%%%%%%%%%%%%%%%%%%%%%%%%%%%%%%%%%%%%%%%%%%%%%%%%%%%
\begin{frame}{Unification}
% [Benedict]

can leave this out...

\end{frame}
%%%%%%%%%%%%%%%%%%%%%%%%%%%%%%%%%%%%%%%%%%%%%%%%%%%%%%%%%%%%%%%%%%%%%%%%%%%%%%%%%%%%%%%%%%


%%%%%%%%%%%%%%%%%%%%%%%%%%%%%%%%%%%%%%%%%%%%%%%%%%%%%%%%%%%%%%%%%%%%%%%%%%%%%%%%%%%%%%%%%%
\begin{frame}{Future plans / wish list}
% [ALL]

\begin{wideitemize}
 \item Unified initial data, source terms, compressible Euler equations (Cartesian)

 \item In-situ data analysis / statistics (similar to \texttt{TimeSeriesCallback} but pointwise for whole field ) \\
       (used to evaluate Held-Suarez type simulations)

 \item In-situ visualization: also for 3D data, requires fast upscaling

 \item Unified post processing: e.g. to visualize projects of the globe

 \item IMEX methods: e.g. Rosenbrock and Multirate Infinitesimal Step (MIS) methods

 \item HEALPix grids
 \begin{itemize}
  \item alternative to cubed sphere geometry
  \item HERA5: ERA5 data on HEALPix grid
 \end{itemize}
\end{wideitemize}

\end{frame}
%%%%%%%%%%%%%%%%%%%%%%%%%%%%%%%%%%%%%%%%%%%%%%%%%%%%%%%%%%%%%%%%%%%%%%%%%%%%%%%%%%%%%%%%%%


%%%%%%%%%%%%%%%%%%%%%%%%%%%%%%%%%%%%%%%%%%%%%%%%%%%%%%%%%%%%%%%%%%%%%%%%%%%%%%%%%%%%%%%%%%
\begin{frame}{Session: Quads to triangles}
% [Philipp]


\end{frame}
%%%%%%%%%%%%%%%%%%%%%%%%%%%%%%%%%%%%%%%%%%%%%%%%%%%%%%%%%%%%%%%%%%%%%%%%%%%%%%%%%%%%%%%%%%



%%%%%%%%%%%%%%%%%%%%%%%%%%%%%%%%%%%%%%%%%%%%%%%%%%%%%%%%%%%%%%%%%%%%%%%%%%%%%%%%%%%%%%%%%%
\begin{frame}{Session: SWE with bottom topography}
% [Tristan]

- try Andres topography

- if this fail, fall back to define your own mountain

\end{frame}
%%%%%%%%%%%%%%%%%%%%%%%%%%%%%%%%%%%%%%%%%%%%%%%%%%%%%%%%%%%%%%%%%%%%%%%%%%%%%%%%%%%%%%%%%%



%%%%%%%%%%%%%%%%%%%%%%%%%%%%%%%%%%%%%%%%%%%%%%%%%%%%%%%%%%%%%%%%%%%%%%%%%%%%%%%%%%%%%%%%%%
\begin{frame}{Session: moist test case [Benedict]}
% [Benedict]

- moist mountain flow

- fallback: dry / wet bubble

\end{frame}
%%%%%%%%%%%%%%%%%%%%%%%%%%%%%%%%%%%%%%%%%%%%%%%%%%%%%%%%%%%%%%%%%%%%%%%%%%%%%%%%%%%%%%%%%%



%%%%%%%%%%%%%%%%%%%%%%%%%%%%%%%%%%%%%%%%%%%%%%%%%%%%%%%%%%%%%%%%%%%%%%%%%%%%%%%%%%%%%%%%%%
\begin{frame}{Session: Vector invariant form}
% [Marco]

- try Andres topography

- if this fail, fall back to define your own mountain

\end{frame}
%%%%%%%%%%%%%%%%%%%%%%%%%%%%%%%%%%%%%%%%%%%%%%%%%%%%%%%%%%%%%%%%%%%%%%%%%%%%%%%%%%%%%%%%%%



%%%%%%%%%%%%%%%%%%%%%%%%%%%%%%%%%%%%%%%%%%%%%%%%%%%%%%%%%%%%%%%%%%%%%%%%%%%%%%%%%%%%%%%%%%
\begin{frame}{Session: Baroclinic instability on steroids}
% [Marco]

- in less than 13 minutes? Woot?

\end{frame}
%%%%%%%%%%%%%%%%%%%%%%%%%%%%%%%%%%%%%%%%%%%%%%%%%%%%%%%%%%%%%%%%%%%%%%%%%%%%%%%%%%%%%%%%%%

\end{document}
